\chapter{Multivariate Distributions}\label{ch:multivariate-distributions}

\chapterinfo{STAT2003, STAT1301}{\cref{ch:probability-spaces,ch:random-variables,ch:expectation-moments,ch:generating-functions,ch:common-distributions}.}

% Sources: STAT2003 handwritten notes (pp. 32-35, 36-39, 40, 46-47); STAT2003 course notes,
% Sections 4.1-4.2; STAT1301 course notes (2026), Sections 3.1-3.4.

In practice, most random experiments involve several random variables. For example:
\begin{enumerate}
  \item select $n = 10$ people at random and observe their heights $X_1, \ldots, X_n$;
  \item toss a coin repeatedly, with $X_i = 1$ if the $i$th toss is Heads and $X_i = 0$ otherwise, giving a sequence $X_1, X_2, \ldots$ of Bernoulli random variables;
  \item select a person at random from a large population and measure their mass $X$ and height $Y$;
  \item simulate 10\,000 realisations $X_1, \ldots, X_{10\,000}$ from the standard normal distribution.
\end{enumerate}
To describe such experiments, it is not enough to specify the distribution of each random variable: we must also say something about how they interact. In the third experiment, if the height $Y$ is large, the mass $X$ is most likely large as well. In the first two, it is reasonable to assume that the random variables are independent, in the sense that information about one gives no information about the others. What we need to specify is the \emph{joint distribution}. The theory follows the same path as for a single random variable.

% ==========================================================================
\section{Joint distributions}\label{sec:joint-distributions}
% ==========================================================================

Let $X_1, \ldots, X_n$ be random variables describing a random experiment. They can be collected into a \emph{random vector} $\vect{X} = (X_1, \ldots, X_n)\T$, a function $\vect{X} \colon \Omega \to \R^n$.

\begin{definition}[Joint CDF]\label{def:joint-cdf}
The \emph{joint cumulative distribution function} of $X_1, \ldots, X_n$ is
\[
  F(x_1, \ldots, x_n) = \Prob(\set{X_1 \le x_1} \cap \cdots \cap \set{X_n \le x_n}) = \Prob(X_1 \le x_1, \ldots, X_n \le x_n).
\]
\end{definition}

The right-hand side shows the abbreviation, used from now on, of commas for intersections. If we know $F$, we can in principle derive any probability involving the $X_i$. As in the one-dimensional case, we distinguish discrete and continuous joint distributions.

\subsection{Discrete joint distributions}\label{sec:discrete-joint}

\begin{definition}[Joint pmf]\label{def:joint-pmf}
The \emph{joint probability mass function} of discrete random variables $X_1, \ldots, X_n$ is
\[
  f(x_1, \ldots, x_n) = \Prob(X_1 = x_1, \ldots, X_n = x_n).
\]
We write $f_{X_1, \ldots, X_n}$, or $f_{\vect{X}}$ for the random vector $\vect{X}$, when the random variables need naming.
\end{definition}

For two random variables $X$ and $Y$,
\[
  f_{X,Y}(x, y) = \Prob(X = x, Y = y) = \Prob\bigl(\set{\omega : X(\omega) = x} \cap \set{\omega : Y(\omega) = y}\bigr),
\]
and $\sum_x \sum_y f_{X,Y}(x, y) = 1$.
By the sum rule, the joint pmf gives every probability involving the random variables. In two dimensions,
\begin{equation}\label{eq:joint-pmf-sum-rule}
  \Prob\bigl((X, Y) \in B\bigr) = \sum_{(x,y) \in B} \Prob(X = x, Y = y)
\end{equation}
for any set $B$ of possible values of $(X, Y)$.

\begin{definition}[Marginal distribution]\label{def:marginal}
The distributions of the individual random variables $X_1, \ldots, X_n$ are their \emph{marginal} distributions. The marginal pmf of $X$ is found by summing the joint pmf over all values of the other variables:
\[
  f_X(x) = \Prob(X = x) = \sum_y \Prob(X = x, Y = y) = \sum_y f_{X,Y}(x, y).
\]
\end{definition}

The converse is not true: from the marginal distributions we cannot in general reconstruct the joint distribution, because we are missing the information on the dependence between the variables. The exception is when the random variables are independent (\cref{sec:independent-rvs}).

\begin{example}[A joint pmf]\label{ex:joint-pmf-table}
Let $X$ take the values $1, 2, 3$ and $Y$ the values $1, 2$, with joint pmf
\begin{center}
\begin{tabular}{@{}c|ccc|c@{}}
 & $x = 1$ & $x = 2$ & $x = 3$ & $f_Y(y)$ \\
\hline
$y = 1$ & 0.1 & 0.3 & 0.2 & 0.6 \\
$y = 2$ & 0.2 & 0.1 & 0.1 & 0.4 \\
\hline
$f_X(x)$ & 0.3 & 0.4 & 0.3 & 1 \\
\end{tabular}
\end{center}
The marginal pmfs are the row and column totals: $f_X = (0.3, 0.4, 0.3)$ on $\set{1,2,3}$ and $f_Y = (0.6, 0.4)$ on $\set{1, 2}$.
\end{example}

\begin{example}[Box with different dice]\label{ex:box-different-dice}
A box contains three dice. Die 1 is an ordinary die; die 2 has no 6 face, but two 5 faces; die 3 has no 5 face, but two 6 faces. A die is selected at random and tossed. Let $X$ be the number of the selected die and $Y$ its face value. The joint pmf, in units of $\tfrac{1}{18}$, is
\begin{center}
\begin{tabular}{@{}c|cccccc|c@{}}
 & $y = 1$ & 2 & 3 & 4 & 5 & 6 & $\sum$ \\
\hline
$x = 1$ & 1 & 1 & 1 & 1 & 1 & 1 & 6 \\
$x = 2$ & 1 & 1 & 1 & 1 & 2 & 0 & 6 \\
$x = 3$ & 1 & 1 & 1 & 1 & 0 & 2 & 6 \\
\hline
$\sum$ & 3 & 3 & 3 & 3 & 3 & 3 & 18 \\
\end{tabular}
\end{center}
So $X$ is uniform on $\set{1, 2, 3}$ and $Y$ is uniform on $\set{1, \ldots, 6}$. It is more convenient to specify the joint probabilities $\Prob(X = x, Y = y)$ than the joint CDF, which follows from them by the sum rule. For example,
\[
  \Prob(X \le 2, Y \le 3) = \Prob(X = 1, Y = 1) + \cdots + \Prob(X = 2, Y = 3) = \tfrac{6}{18} = \tfrac13.
\]
We cannot reconstruct the inside of this table from the row and column totals alone.
\end{example}

\subsection{Continuous joint distributions}\label{sec:continuous-joint}

\begin{definition}[Joint pdf]\label{def:joint-pdf}
Continuous random variables $X_1, \ldots, X_n$ have \emph{joint probability density function} $f$ if
\[
  \Prob(a_1 \le X_1 \le b_1, \ldots, a_n \le X_n \le b_n) = \int_{a_1}^{b_1} \cdots \int_{a_n}^{b_n} f(x_1, \ldots, x_n) \dd x_n \cdots \dd x_1
\]
for all $a_1, \ldots, a_n, b_1, \ldots, b_n$. The total integral of $f$ over $\R^n$ is 1. We write $f_{X_1, \ldots, X_n}$ or $f_{\vect{X}}$ when needed.
\end{definition}

The joint pdf is a continuous analogue of the joint pmf: $f(x_1, \ldots, x_n)$ is the ``density'' that $X_1 = x_1, \ldots, X_n = x_n$. In two dimensions, for small $h$,
\[
  \Prob(x \le X \le x + h, \, y \le Y \le y + h) = \int_x^{x+h} \int_y^{y+h} f(u, v) \dd v \dd u \approx h^2 f(x, y).
\]
Given the joint pdf, every probability can be calculated by integration. In two dimensions,
\begin{equation}\label{eq:joint-pdf-integral-rule}
  \Prob\bigl((X, Y) \in B\bigr) = \iint_{B} f(x, y) \dd x \dd y
\end{equation}
for any subset $B$ of $\R^2$. Instead of the area under a curve, we now evaluate a volume under a surface. As in the discrete case, the marginal pdf of each $X_i$ is found by integrating $f$ over all the other variables; in two dimensions,
\[
  f_X(x) = \int_{-\infty}^\infty f(x, y) \dd y.
\]
Again, we usually cannot reconstruct the joint pdf from the marginal pdfs unless the random variables are independent.

% ==========================================================================
\section{Independent random variables}\label{sec:independent-rvs}
% ==========================================================================

Random variables are independent if information about one of them does not affect our knowledge of the others.

\begin{definition}[Independent random variables]\label{def:independent-rvs}
Random variables $X_1, \ldots, X_n$ are \emph{independent} if, for all $A_1, \ldots, A_n \subseteq \R$,
\[
  \Prob(X_1 \in A_1, \ldots, X_n \in A_n) = \Prob(X_1 \in A_1) \cdots \Prob(X_n \in A_n).
\]
That is, the events $\set{X_1 \in A_1}, \ldots, \set{X_n \in A_n}$ are independent for any choice of the sets $A_i$. An infinite sequence $X_1, X_2, \ldots$ is independent if, for any finite choice of distinct indices $i_1, \ldots, i_n$, the random variables $X_{i_1}, \ldots, X_{i_n}$ are independent.
\end{definition}

\begin{theorem}[Independent discrete random variables]\label{thm:independent-discrete}
Discrete random variables $X_1, \ldots, X_n$ are independent if and only if
\begin{equation}\label{eq:independence-pmf}
  \Prob(X_1 = x_1, \ldots, X_n = x_n) = \Prob(X_1 = x_1) \cdots \Prob(X_n = x_n)
\end{equation}
for all $x_1, \ldots, x_n$. That is, $f_{X_1, \ldots, X_n}(x_1, \ldots, x_n) = f_{X_1}(x_1) \cdots f_{X_n}(x_n)$.
\end{theorem}

\begin{proof}
If $X_1, \ldots, X_n$ are independent, then the events $\set{X_1 = x_1}, \ldots, \set{X_n = x_n}$ are independent, which gives \eqref{eq:independence-pmf}. Conversely, write
\[
  \Prob(X_1 \in A_1, \ldots, X_n \in A_n) = \sum_{x_1 \in A_1} \cdots \sum_{x_n \in A_n} \Prob(X_1 = x_1, \ldots, X_n = x_n).
\]
If \eqref{eq:independence-pmf} holds, the multiple sum factorises as
\[
  \sum_{x_1 \in A_1} \Prob(X_1 = x_1) \cdots \sum_{x_n \in A_n} \Prob(X_n = x_n) = \Prob(X_1 \in A_1) \cdots \Prob(X_n \in A_n). \qedhere
\]
\end{proof}

\begin{theorem}[Independent continuous random variables]\label{thm:independent-continuous}
Let $X_1, \ldots, X_n$ be continuous random variables with joint pdf $f$ and marginal pdfs $f_{X_1}, \ldots, f_{X_n}$. Then $X_1, \ldots, X_n$ are independent if and only if
\[
  f(x_1, \ldots, x_n) = f_{X_1}(x_1) \cdots f_{X_n}(x_n) \qquad \text{for all } x_1, \ldots, x_n.
\]
\end{theorem}

\begin{example}[Box with ordinary dice]\label{ex:box-same-dice}
Repeat \cref{ex:box-different-dice} with three ordinary fair dice. Now the events $\set{X = x}$ and $\set{Y = y}$ are independent, so each entry of the joint pmf table is $\tfrac13 \times \tfrac16$. In the original experiment, not all of these events are independent.
\end{example}

\begin{example}[Coin flips]\label{ex:coin-flips-independent}
Flip a coin $n$ times. For $i = 1, \ldots, n$, let $X_i = 1$ if the $i$th toss is Heads and $X_i = 0$ if Tails, with $\Prob(X_i = 1) = p = 1 - \Prob(X_i = 0)$, and assume $X_1, \ldots, X_n$ are independent. This completely specifies the model. For example, let $X = X_1 + \cdots + X_n$ be the total number of Heads, and let $A$ be the set of binary vectors $\vect{x} = (x_1, \ldots, x_n)$ with $\sum_i x_i = k$, which has $\binom{n}{k}$ elements. Then
\begin{align*}
  \Prob(X = k) &= \sum_{\vect{x} \in A} \Prob(X_1 = x_1, \ldots, X_n = x_n) = \sum_{\vect{x} \in A} \Prob(X_1 = x_1) \cdots \Prob(X_n = x_n) \\
  &= \sum_{\vect{x} \in A} p^k (1-p)^{n-k} = \binom{n}{k} p^k (1-p)^{n-k}.
\end{align*}
So $X \sim \Bin(n, p)$ (compare \cref{ex:binomial-law}).
\end{example}

Many statistical models involve random variables $X_1, X_2, \ldots$ that are independent and identically distributed (iid, \cref{def:iid}). We write the model as
\[
  X_1, X_2, \ldots \iid \text{Dist} \qquad \text{or} \qquad X_1, X_2, \ldots \iid f \qquad \text{or} \qquad X_1, X_2, \ldots \iid F,
\]
where Dist is the common distribution, with pdf $f$ and CDF $F$.

\begin{remark}[Random sample]\label{rem:random-sample}
The model in which $X_1, \ldots, X_n$ are independent and all have the same distribution is the most widely used model in statistics. We say $X_1, \ldots, X_n$ is a \emph{random sample} of size $n$ from the given distribution, which is sometimes called the \emph{sampling distribution}. In \cref{ex:coin-flips-independent}, $X_1, \ldots, X_n$ is a random sample from the $\Bern(p)$ distribution.
\end{remark}

\begin{example}[Uniform numbers]\label{ex:uniform-numbers}
Select $n$ points randomly and independently from $[0, 1]$, for example with a computer's random number generator. A model is: $X_1, \ldots, X_n$ are independent, each $\Unif[0, 1]$. The joint pdf is simply
\[
  f(x_1, \ldots, x_n) = 1, \qquad 0 \le x_1 \le 1, \ldots, 0 \le x_n \le 1
\]
(and 0 elsewhere). In principle we can now calculate any probability involving the $X_i$. For example, for $n = 2$, by \eqref{eq:joint-pdf-integral-rule},
\[
  \Prob\Bigl(\frac{X_1 + X_2^2}{X_1 X_2} > \sin(X_1 - X_2)\Bigr) = \iint_A 1 \dd x_1 \dd x_2 = \operatorname{Area}(A),
\]
where
\[
  A = \Bigl\{(x_1, x_2) \in [0, 1]^2 : \frac{x_1 + x_2^2}{x_1 x_2} > \sin(x_1 - x_2)\Bigr\}.
\]
On a computer we can generate random samples from many distributions. The R code below draws two random samples of size 1000, one from $\Unif[0, 1]$ and one from $\Normal(\tfrac12, \tfrac{1}{12})$, and plots their histograms with 20 bins. The two distributions have the same expectation and variance, but the ``density'' of points in the two samples is clearly different, and follows the shape of the corresponding pdfs.
\begin{rcode}
x <- runif(1000)                         # U[0, 1] sample
hist(x, breaks = 20)
y <- 0.5 + rnorm(1000) * sqrt(1/12)      # N(1/2, 1/12) sample
hist(y, breaks = 20)
\end{rcode}
\end{example}

\begin{example}[Bivariate standard normal distribution]\label{ex:bivariate-standard-normal}
Suppose $X$ and $Y$ are independent, both with a standard normal distribution. Then $(X, Y)$ has a \emph{bivariate standard normal distribution}, with joint pdf
\[
  f(x, y) = f_X(x) f_Y(y) = \frac{1}{\sqrt{2\pi}} e^{-x^2/2} \cdot \frac{1}{\sqrt{2\pi}} e^{-y^2/2} = \frac{1}{2\pi} e^{-(x^2 + y^2)/2}.
\]
Its graph is a hat-shaped surface. The contour lines of the joint pdf are circles, and a scatter plot of simulated pairs shows a ``spherical'' pattern:
\begin{rcode}
x <- rnorm(2000)
y <- rnorm(2000)
plot(y ~ x, xlim = c(-3, 3), ylim = c(-3, 3))
\end{rcode}
\end{example}

% ==========================================================================
\section{Expectation}\label{sec:joint-expectation}
% ==========================================================================

As in the one-dimensional case, the expected value of a real-valued function of $X_1, \ldots, X_n$ is the weighted average of all the values the function can take. If $Z = h(X_1, \ldots, X_n)$, then
\[
  \E Z = \sum_{x_1} \cdots \sum_{x_n} h(x_1, \ldots, x_n) f(x_1, \ldots, x_n)
\]
in the discrete case, where $f$ is the joint pmf, and
\[
  \E Z = \int_{-\infty}^\infty \cdots \int_{-\infty}^\infty h(x_1, \ldots, x_n) f(x_1, \ldots, x_n) \dd x_1 \cdots \dd x_n
\]
in the continuous case, where $f$ is the joint pdf. Here $h \colon \R^n \to \R$: the random vector maps $\Omega$ into $\R^n$, and $h$ maps $\R^n$ into $\R$.

\begin{theorem}[Properties of the expectation]\label{thm:joint-expectation-properties}
Let $X_1, \ldots, X_n$ be discrete or continuous random variables with expectations $\mu_1, \ldots, \mu_n$, and let $Y = a + b_1 X_1 + \cdots + b_n X_n$ for constants $a, b_1, \ldots, b_n$. Then
\[
  \E Y = a + b_1 \E X_1 + \cdots + b_n \E X_n = a + b_1 \mu_1 + \cdots + b_n \mu_n.
\]
In particular, $\E(X + Y) = \E X + \E Y$. This holds whether or not the random variables are independent. If $X_1, \ldots, X_n$ are independent, then also
\[
  \E X_1 X_2 \cdots X_n = \E X_1 \, \E X_2 \cdots \E X_n.
\]
\end{theorem}

\begin{proof}
For the sum of two discrete random variables with joint pmf $f$,
\begin{align*}
  \E(X + Y) &= \sum_x \sum_y (x + y) f_{X,Y}(x, y) = \sum_x \sum_y x f_{X,Y}(x, y) + \sum_x \sum_y y f_{X,Y}(x, y) \\
  &= \sum_x x \sum_y f_{X,Y}(x, y) + \sum_y y \sum_x f_{X,Y}(x, y) = \sum_x x f_X(x) + \sum_y y f_Y(y) = \E X + \E Y.
\end{align*}
The continuous case is the same with integrals, and the general linear case follows in the same way. For the product, take two independent continuous random variables, so that $f(x, y) = f_X(x) f_Y(y)$. Then
\[
  \E XY = \iint x y f(x, y) \dd x \dd y = \iint x y f_X(x) f_Y(y) \dd x \dd y = \int x f_X(x) \dd x \int y f_Y(y) \dd y = \E X \, \E Y.
\]
The generalisation to $n$ variables, and to the discrete case, is obvious.
\end{proof}

\begin{example}[Expectation of the binomial and hypergeometric distributions]\label{ex:binomial-hypergeometric-expectation}
If $X \sim \Bin(n, p)$, then $X$ is a sum $X_1 + \cdots + X_n$ of iid $\Bern(p)$ random variables (\cref{ex:coin-flips-independent}). So
\[
  \E X = \E X_1 + \cdots + \E X_n = p + \cdots + p = np.
\]
This did not use the independence of the $X_i$. Now let $X \sim \Hyp(n, r, N)$ and $p = r/N$: $X$ is the number of red balls when $n$ balls are drawn from an urn with $r$ red and $N - r$ other balls. Without loss of generality the balls are drawn one by one; let $X_i = 1$ if the $i$th ball is red and 0 otherwise. Then again $X = X_1 + \cdots + X_n$ with each $X_i \sim \Bern(p)$, but now the $X_i$ are dependent. This does not affect the calculation, so $\E X = np = nr/N$.
\end{example}

The product rule for independent random variables is particularly handy in combination with transforms.

\begin{example}[Sum of Poisson random variables]\label{ex:sum-poisson-pgf}
Let $X \sim \Pois(\lambda)$ and $Y \sim \Pois(\mu)$ be independent. Their PGFs are $e^{-\lambda(1-z)}$ and $e^{-\mu(1-z)}$ (\cref{ex:pgf-poisson}), so the PGF of $X + Y$ is
\[
  \E z^{X+Y} = \E z^X \, \E z^Y = e^{-\lambda(1-z)} e^{-\mu(1-z)} = e^{-(\lambda + \mu)(1-z)}.
\]
By the uniqueness property of the PGF, $X + Y \sim \Pois(\lambda + \mu)$. In general, if $X$ and $Y$ are independent, then $G_{X+Y}(z) = G_X(z) G_Y(z)$ and $M_{X+Y}(s) = M_X(s) M_Y(s)$.
\end{example}

\begin{example}[Sum of exponential random variables]\label{ex:sum-exponential-mgf}
The MGF of $X \sim \GammaDist(\alpha, \lambda)$ is $\bigl(\lambda/(\lambda - s)\bigr)^\alpha$ (\cref{thm:gamma-properties}), so the MGF of the $\Expo(\lambda)$ distribution is $\lambda/(\lambda - s)$. Let $X_1, \ldots, X_n$ be iid $\Expo(\lambda)$ and $S_n = X_1 + \cdots + X_n$. Then
\[
  \E e^{sS_n} = \E\bigl[e^{sX_1} \cdots e^{sX_n}\bigr] = \E e^{sX_1} \cdots \E e^{sX_n} = \Bigl(\frac{\lambda}{\lambda - s}\Bigr)^n,
\]
so $S_n \sim \GammaDist(n, \lambda)$.
\end{example}

% ==========================================================================
\section{Covariance and correlation}\label{sec:covariance}
% ==========================================================================

\begin{definition}[Covariance]\label{def:covariance}
The \emph{covariance} of random variables $X$ and $Y$, with $\E X = \mu_X$ and $\E Y = \mu_Y$, is
\[
  \Cov(X, Y) \coloneqq \E\bigl[(X - \mu_X)(Y - \mu_Y)\bigr].
\]
\end{definition}

The covariance measures the amount of \emph{linear} dependence between the variables. If small values of $X$ (below its expectation) tend to go with small values of $Y$, and large values of $X$ with large values of $Y$, then $\Cov(X, Y)$ is positive. If small values of $X$ go with large values of $Y$, and large with small, then $\Cov(X, Y)$ is negative.

\begin{theorem}[Properties of variance and covariance]\label{thm:covariance-properties}
For random variables $X$, $Y$ and $Z$ and constants $a$ and $b$:
\begin{enumerate}
  \item $\Var(X) = \E X^2 - (\E X)^2$;
  \item $\Var(aX + b) = a^2 \Var(X)$;
  \item $\Cov(X, Y) = \E XY - \E X \, \E Y$;
  \item $\Cov(X, Y) = \Cov(Y, X)$;
  \item $\Cov(aX + bY, Z) = a \Cov(X, Z) + b \Cov(Y, Z)$;
  \item $\Cov(X, X) = \Var(X)$;
  \item $\Var(X + Y) = \Var(X) + \Var(Y) + 2\Cov(X, Y)$;
  \item if $X$ and $Y$ are independent, then $\Cov(X, Y) = 0$.
\end{enumerate}
\end{theorem}

\begin{proof}
Properties 1 and 2 are \cref{thm:variance-properties}. Write $\mu_X = \E X$ and $\mu_Y = \E Y$.

3. $\Cov(X, Y) = \E[XY - Y\mu_X - X\mu_Y + \mu_X\mu_Y] = \E XY - \mu_Y\mu_X - \mu_X\mu_Y + \mu_X\mu_Y = \E XY - \mu_X\mu_Y$.

4. $\Cov(X, Y) = \E[(X - \mu_X)(Y - \mu_Y)] = \E[(Y - \mu_Y)(X - \mu_X)] = \Cov(Y, X)$.

5. Using property 3,
\begin{align*}
  \Cov(aX + bY, Z) &= \E[(aX + bY)Z] - \E[aX + bY]\,\E Z \\
  &= a\E XZ + b \E YZ - a \E X \, \E Z - b \E Y \, \E Z \\
  &= a\bigl(\E XZ - \E X\,\E Z\bigr) + b\bigl(\E YZ - \E Y\,\E Z\bigr) = a\Cov(X, Z) + b\Cov(Y, Z).
\end{align*}

6. $\Cov(X, X) = \E[(X - \mu_X)^2] = \Var(X)$.

7. By property 6, $\Var(X + Y) = \Cov(X + Y, X + Y)$. By property 5, and property 4 to expand the second argument,
\[
  \Cov(X + Y, X + Y) = \Cov(X, X) + \Cov(X, Y) + \Cov(Y, X) + \Cov(Y, Y),
\]
which by properties 4 and 6 equals $\Var(X) + \Var(Y) + 2\Cov(X, Y)$.

8. If $X$ and $Y$ are independent, then $\E XY = \E X\,\E Y$ (\cref{thm:joint-expectation-properties}), so $\Cov(X, Y) = 0$ by property 3.
\end{proof}

A scaled version of the covariance is the correlation coefficient.

\begin{definition}[Correlation coefficient]\label{def:correlation}
The \emph{correlation coefficient} of $X$ and $Y$, with standard deviations $\sigma_X$ and $\sigma_Y$, is
\[
  \Corr(X, Y) \coloneqq \frac{\Cov(X, Y)}{\sqrt{\Var(X)}\sqrt{\Var(Y)}} = \frac{\Cov(X, Y)}{\sigma_X \sigma_Y}.
\]
\end{definition}

\begin{theorem}\label{thm:correlation-bounds}
The correlation coefficient always lies between $-1$ and $1$.
\end{theorem}

\begin{proof}
Let $a > 0$. The variance of $\pm aX + Y$ is non-negative, so by the properties of variance and covariance,
\[
  \Var(-aX + Y) = a^2\sigma_X^2 + \sigma_Y^2 - 2a\Cov(X, Y) \ge 0.
\]
Rearranging and dividing by $2a\sigma_X\sigma_Y$ gives
\[
  \Corr(X, Y) \le \frac12\Bigl(\frac{a\sigma_X}{\sigma_Y} + \frac{\sigma_Y}{a\sigma_X}\Bigr).
\]
Similarly, $\Var(aX + Y) = a^2\sigma_X^2 + \sigma_Y^2 + 2a\Cov(X, Y) \ge 0$ gives
\[
  \Corr(X, Y) \ge -\frac12\Bigl(\frac{a\sigma_X}{\sigma_Y} + \frac{\sigma_Y}{a\sigma_X}\Bigr).
\]
Choosing $a = \sigma_Y/\sigma_X$ makes the bracket equal to $2$, so $-1 \le \Corr(X, Y) \le 1$.
\end{proof}

Combining properties 2, 5 and 7 gives the variance of a linear combination.

\begin{theorem}[Variance of a linear combination]\label{thm:variance-linear-combination}
Let $X_1, \ldots, X_n$ be independent random variables with variances $\sigma_1^2, \ldots, \sigma_n^2$, and let $Y = a + b_1 X_1 + \cdots + b_n X_n$ for constants $a, b_1, \ldots, b_n$. Then
\[
  \Var(Y) = b_1^2 \Var(X_1) + \cdots + b_n^2 \Var(X_n) = b_1^2\sigma_1^2 + \cdots + b_n^2\sigma_n^2.
\]
\end{theorem}

\begin{proof}
The constant $a$ plays no role in the variance. By property 6 and repeated application of property 5,
\[
  \Var(Y) = \Cov\Bigl(\sum_i b_i X_i, \sum_j b_j X_j\Bigr) = \sum_{i=1}^n \Cov(b_i X_i, b_i X_i) + 2\sum_{i<j} \Cov(b_i X_i, b_j X_j).
\]
Here $\Cov(b_i X_i, b_i X_i) = b_i^2 \Var(X_i)$, and all the covariance terms with $i \ne j$ are zero by independence.
\end{proof}

\begin{example}[Variance of the binomial and hypergeometric distributions]\label{ex:binomial-hypergeometric-variance}
As in \cref{ex:binomial-hypergeometric-expectation}, write $X = X_1 + \cdots + X_n$ with each $X_i \sim \Bern(p)$; the $X_i$ are independent in the binomial case and dependent in the hypergeometric case, where $p = r/N$. For the binomial distribution, \cref{thm:variance-linear-combination} gives directly
\[
  \Var(X) = \Var(X_1) + \cdots + \Var(X_n) = n\Var(X_1) = np(1-p).
\]
In the hypergeometric case the covariance terms must be included:
\[
  \Var(X) = \Var(X_1) + \cdots + \Var(X_n) + 2\sum_{i<j} \Cov(X_i, X_j).
\]
By symmetry all the $\Cov(X_i, X_j)$, $i \ne j$, are the same, so
\[
  \Var(X) = n\Var(X_1) + n(n-1)\Cov(X_1, X_2).
\]
Here $\Var(X_1) = p(1-p)$ and $\Cov(X_1, X_2) = \E X_1X_2 - p^2$, with
\[
  \E X_1 X_2 = \Prob(X_1 = 1, X_2 = 1) = \Prob(X_1 = 1)\,\Prob(X_2 = 1 \given X_1 = 1) = p\,\frac{r - 1}{N - 1}.
\]
Simplifying gives $\Var(X) = np(1-p)\dfrac{N - n}{N - 1}$.
\end{example}

% ==========================================================================
\section{Expectation vectors and covariance matrices}\label{sec:covariance-matrix}
% ==========================================================================

For a random vector $\vect{X} = (X_1, \ldots, X_n)\T$ it is convenient to write expectations and covariances in vector notation. The expectation of a vector or matrix of random variables is the vector or matrix of their expectations.

\begin{definition}[Expectation vector, covariance matrix]\label{def:covariance-matrix}
The \emph{expectation vector} of a random vector $\vect{X}$ is
\[
  \vect{\mu} = (\mu_1, \ldots, \mu_n)\T \coloneqq (\E X_1, \ldots, \E X_n)\T = \E\vect{X} \in \R^n.
\]
The \emph{covariance matrix} $\mat{\Sigma} = \Cov(\vect{X}) \in \R^{n \times n}$ is the matrix with $(i,j)$th element
\[
  \Sigma_{ij} = \Cov(X_i, X_j) = \E(X_i - \mu_i)(X_j - \mu_j) = \sigma_i\sigma_j\Corr(X_i, X_j),
\]
where $\sigma_i^2 = \Var(X_i)$. In matrix form, $\mat{\Sigma} = \E(\vect{X} - \vect{\mu})(\vect{X} - \vect{\mu})\T$. We write $\vect{\mu}_{\vect{X}}$ and $\mat{\Sigma}_{\vect{X}}$ to emphasise the random vector.
\end{definition}

$\vect{\mu}$ and $\mat{\Sigma}$ play the same roles as $\mu$ and $\sigma^2$ in the one-dimensional case. The covariance matrix has the following properties.

\begin{theorem}[Properties of covariance matrices]\label{thm:covariance-matrix-properties}
Let $\vect{X}, \vect{Y} \in \R^n$ be random vectors and $\vect{b} \in \R^n$ a constant vector.
\begin{enumerate}
  \item $\Cov(\vect{X}) = \E(\vect{X}\vect{X}\T) - (\E\vect{X})(\E\vect{X})\T$. (Compare $\Var(X) = \E X^2 - (\E X)^2$ and $\Cov(X, Y) = \E XY - \E X\,\E Y$.)
  \item $\mat{\Sigma}$ is symmetric: $\mat{\Sigma}\T = \mat{\Sigma}$.
  \item $\mat{\Sigma}$ is positive semidefinite: $\vect{u}\T\mat{\Sigma}\vect{u} \ge 0$ for every vector $\vect{u} \in \R^n$.
  \item If $X_1, \ldots, X_n$ are independent, then $\mat{\Sigma}$ is a diagonal matrix.
  \item $\bigl[\Cov(\vect{X} + \vect{Y})\bigr]_{ij} = \bigl[\Cov(\vect{X})\bigr]_{ij} + \bigl[\Cov(\vect{Y})\bigr]_{ij} + \Cov(X_i, Y_j) + \Cov(Y_i, X_j)$. So $\Cov(\vect{X} + \vect{Y}) = \Cov(\vect{X}) + \Cov(\vect{Y})$ if and only if $X_i$ and $Y_j$ are uncorrelated for all $i$ and $j$.
  \item $\Cov(\vect{X} + \vect{b}) = \Cov(\vect{X})$.
  \item If $\mat{A}$ is an $m \times n$ matrix and $\vect{Y} = \mat{A}\vect{X}$, then $\Cov(\vect{Y}) = \mat{A}\Cov(\vect{X})\mat{A}\T$.
\end{enumerate}
\end{theorem}

\begin{proof}
1. Writing $\mu_i = \E X_i$,
\[
  \E(\vect{X}\vect{X}\T) - (\E\vect{X})(\E\vect{X})\T = \E\begin{pmatrix} X_1X_1 & X_1X_2 & \cdots \\ X_2X_1 & X_2X_2 & \cdots \\ \vdots & & \ddots \end{pmatrix} - \begin{pmatrix} \mu_1\mu_1 & \mu_1\mu_2 & \cdots \\ \mu_2\mu_1 & \mu_2\mu_2 & \cdots \\ \vdots & & \ddots \end{pmatrix},
\]
whose $(i,j)$th element is $\E X_iX_j - \mu_i\mu_j = \Cov(X_i, X_j)$ by \cref{thm:covariance-properties}(3).

2. $\Cov(X_i, X_j) = \Cov(X_j, X_i)$.

3. Let $\vect{Y} = \vect{X} - \vect{\mu}$. Then $\vect{Y}\T\vect{u}$ is a random variable, and
\[
  \vect{u}\T\mat{\Sigma}\vect{u} = \vect{u}\T\E(\vect{Y}\vect{Y}\T)\vect{u} = \E(\vect{u}\T\vect{Y}\vect{Y}\T\vect{u}) = \E\bigl((\vect{Y}\T\vect{u})\T(\vect{Y}\T\vect{u})\bigr) = \E(\vect{Y}\T\vect{u})^2 \ge 0.
\]

4. Independent random variables have zero covariance (\cref{thm:covariance-properties}(8)).

5. Expand $\Cov(X_i + Y_i, X_j + Y_j)$ using \cref{thm:covariance-properties}(5) and (4).

6. Adding a constant does not change $\vect{X} - \E\vect{X}$.

7. By linearity, $\E\vect{Y} = \mat{A}\E\vect{X}$, so
\begin{align*}
  \Cov(\vect{Y}) &= \E\bigl[(\vect{Y} - \E\vect{Y})(\vect{Y} - \E\vect{Y})\T\bigr] = \E\bigl[(\mat{A}\vect{X} - \mat{A}\E\vect{X})(\mat{A}\vect{X} - \mat{A}\E\vect{X})\T\bigr] \\
  &= \E\bigl[\mat{A}(\vect{X} - \E\vect{X})(\vect{X} - \E\vect{X})\T\mat{A}\T\bigr] \\
  &= \mat{A}\,\E\bigl[(\vect{X} - \E\vect{X})(\vect{X} - \E\vect{X})\T\bigr]\mat{A}\T = \mat{A}\Cov(\vect{X})\mat{A}\T. \qedhere
\end{align*}
\end{proof}

\begin{remark}
Since $\mat{\Sigma}$ is symmetric and positive semidefinite, its eigenvalues are real and non-negative.
\end{remark}

% ==========================================================================
\section{The multinomial distribution}\label{sec:multinomial}
% ==========================================================================

An important discrete joint distribution is the multinomial distribution, a generalisation of the binomial distribution.

\begin{definition}[Multinomial distribution]\label{def:multinomial}
$(X_1, \ldots, X_k)$ has a \emph{multinomial distribution} with parameters $n$ and $p_1, \ldots, p_k$, written $(X_1, \ldots, X_k) \sim \Mnom(n, p_1, \ldots, p_k)$, if
\[
  \Prob(X_1 = x_1, \ldots, X_k = x_k) = \frac{n!}{x_1!\,x_2! \cdots x_k!}\,p_1^{x_1} p_2^{x_2} \cdots p_k^{x_k}
\]
for all $x_1, \ldots, x_k \in \set{0, 1, \ldots, n}$ with $x_1 + \cdots + x_k = n$.
\end{definition}

\begin{example}[Balls in urns]\label{ex:multinomial-urns}
Throw $n$ balls independently into $k$ urns, so that each ball lands in urn $i$ with probability $p_i$. Let $X_i$ be the number of balls in urn $i$. Then $(X_1, \ldots, X_k) \sim \Mnom(n, p_1, \ldots, p_k)$. Indeed, let $x_1, \ldots, x_k$ be integers between 0 and $n$ summing to $n$. The probability that the first $x_1$ balls fall in urn 1, the next $x_2$ in urn 2, and so on, is $p_1^{x_1} p_2^{x_2} \cdots p_k^{x_k}$. To get the probability of $x_1$ balls in urn 1, $x_2$ in urn 2, and so on, multiply this by the number of ways to fill the urns with $x_1, \ldots, x_k$ balls, which is $n!/(x_1! \cdots x_k!)$.
\end{example}

\begin{remark}
The binomial distribution is the case of only two urns. For each $i$, $X_i \sim \Bin(n, p_i)$.
\end{remark}
